\section*{第 \ref{ch:b1:reals} 章 --- 实数}

\begin{solution}{exo:b1:reals:1}
$A$：每个元素都 $\leq 1$，且 $1 \in A$：$\sup A = \max A = 1$。
下界：$0$ 为下界；对 $\varepsilon > 0$，由阿基米德
性质给出 $n$ 使 $\frac 1n < \varepsilon$，故没有正数是
$A$ 的下界：$\inf A = 0$，且不能取到（即无最小值）。

$B$：各项为 $0, -\frac12, \frac23, -\frac34, \frac45, \dots$。其中偶数项
$\frac{n}{n+1}$（$n$ 为偶数）递增趋于 $1$，但取不到它；
奇数项 $-\frac{n}{n+1}$ 递减趋于 $-1$。于是 $\sup B = 1$ 且
$\inf B = -1$，二者均取不到，故既无最大值也无最小值。（界：$\abs{b}
< 1$ 对所有 $b \in B$ 成立；且 $\frac{n}{n+1} = 1 - \frac{1}{n+1} > 1 -
\varepsilon$ 对充分大的 $n$ 成立，下界同理。）

$C = \intoo{-\sqrt 3}{\sqrt 3}$：$\sup C = \sqrt 3$，$\inf C =
-\sqrt 3$，二者均取不到。
\end{solution}

\begin{solution}{exo:b1:reals:2}
记 $x = \lfloor x \rfloor + u$，$y = \lfloor y \rfloor + v$，
其中 $u, v \in \intco{0}{1}$。于是 $x + y = \lfloor x \rfloor + \lfloor y
\rfloor + (u + v)$，其中 $u + v \in \intco{0}{2}$。若 $u + v < 1$，
则 $\lfloor x + y\rfloor = \lfloor x\rfloor + \lfloor y \rfloor$；若
$1 \leq u + v < 2$，则 $\lfloor x+y \rfloor = \lfloor x \rfloor +
\lfloor y \rfloor + 1$。两种情形均可出现：$(x, y) = (0.2,\, 0.3)$
给出左侧等式，$(0.7,\, 0.8)$ 给出右侧者。
\end{solution}

\begin{solution}{exo:b1:reals:3}
设 $k = \lfloor x \rfloor$，则 $k \leq x < k + 1$。两边乘以
$n$：$nk \leq nx < nk + n$，再取整
（取整作用于整数一侧时保持不等号），得 $nk \leq \lfloor nx \rfloor \leq nk + n
- 1$。两边除以 $n$：$k \leq \frac{\lfloor nx \rfloor}{n} < k + 1$，
故外层取整为 $k$。
\end{solution}

\begin{solution}{exo:b1:reals:4}
$A$ 的每个元素都属于 $B$，因此 $\sup B$ 为 $A$ 的一个上界：
从而 $\sup A \leq \sup B$（$\sup A$ 是\emph{最小上界}
）。同理 $\inf B \leq \inf A$。最后 $\inf A \leq
\sup A$，因 $A$ 非空：任一 $a \in A$ 都位于二者之间。
\end{solution}

\begin{solution}{exo:b1:reals:5}
设 $s = \sup A$，$t = \sup B$。每个和 $a + b \leq s + t$，即为
上界。对 $\varepsilon > 0$，取 $a > s - \frac\varepsilon2$ 与
$b > t - \frac\varepsilon2$（见\cref{prop:b1:reals:epsilon}），则 $a +
b > s + t - \varepsilon$。由上确界的 $\varepsilon$ 刻画，
$\sup(A+B) = s + t$。

对 $-A$：$m$ 为 $-A$ 的上界 $\iff$ $-m$ 为 $A$ 的下界；
故 $-A$ 的最小上界对应
$A$ 的最大下界，即 $\sup(-A) = -\inf A$。
\end{solution}

\begin{solution}{exo:b1:reals:6}
对每个 $x$：$f(x) + g(x) \leq \sup f + \sup g$；对左端取
上确界即得该不等式。取一个严格不等式的例子：$E = \{0, 1\}$，
$f = \mathbf{1}_{\{0\}}$（取值 $1$ 于 $0$ 处，其余处为 $0$），$g =
\mathbf{1}_{\{1\}}$，则 $\sup(f + g) = 1 < 2 = \sup f + \sup g$。

与\cref{exo:b1:reals:5}并不矛盾：在那里，$a \in A$ 与 $b
\in B$ \emph{彼此独立}地变化；这里同一个 $x$ 同时作用于 $f$
与 $g$ --- 集合 $\{f(x) + g(x) : x \in E\}$ 比
集合 $\{f(x) + g(y) : x, y \in E\}$ 更小。
\end{solution}

\begin{solution}{exo:b1:reals:7}
$\sqrt 6$ 是无理数：$6 = 2 \times 3$ 不是完全平方数，
且 $v_2(6q^2) = 1 + 2v_2(q)$ 为奇数，故 $6q^2 = r^2$ 无解
（参见\cref{exo:b1:arith:7}）。设 $x = \sqrt 2 + \sqrt 3 \in \Q$。
于是 $x^2 = 5 + 2\sqrt 6 \in \Q$，故 $\sqrt 6 = \frac{x^2 - 5}{2} \in
\Q$，矛盾。因此 $\sqrt 2 + \sqrt 3 \notin \Q$。
\end{solution}

\begin{solution}{exo:b1:reals:8}
设 $x < y$。取 $n \in \N$ 使 $2^n > \frac{1}{y - x}$（阿基米德性质：
$2^n \geq n + 1$，由归纳法易证，故某个 $2$ 的幂超过
任一实数）。于是，仿照\cref{thm:b1:reals:density}的证明，
以 $2^n$ 代替 $n$：$m = \lfloor 2^n x \rfloor + 1$ 满足 $x <
\frac{m}{2^n} < y$。故 $D$ 稠密。
\end{solution}

\begin{solution}{exo:b1:reals:9}
\begin{enumerate}
  \item 设 $\alpha > 0$。先证 $\alpha \in G$。若不然，由
        下确界的 $\varepsilon$ 刻画，取
        $\varepsilon = \alpha$，存在 $g \in G$ 使 $\alpha < g <
        2\alpha$（左端严格，因 $\alpha \notin G$）；再取
        $\varepsilon = g - \alpha$，存在 $h \in G$ 使
        $\alpha < h < g$。于是 $g - h \in G$，且 $0 < g - h < g -
        \alpha < \alpha$：$G \cap \intoo{0}{+\infty}$ 中的元素
        位于其下确界之下，矛盾。故 $\alpha \in G$，且 $\alpha\Z
        \subseteq G$（$G$ 为群）。反之，对 $x \in G$，令
        $k = \lfloor x/\alpha \rfloor$：则 $x - k\alpha \in G$ 且
        $0 \leq x - k\alpha < \alpha$，由 $\alpha$ 的定义，
        必有 $x - k\alpha = 0$。因此 $G = \alpha\Z$。
  \item 设 $\alpha = 0$，且 $x < y$。存在 $g \in G$
        使 $0 < g < y - x$。倍数 $kg$（$k = \lfloor x/g
        \rfloor + 1$）满足 $x < kg \leq x + g < y$，且 $kg \in
        G$，故稠密。
\end{enumerate}
$G = \Z + \sqrt 2\,\Z$ 是 $(\R, +)$ 的子群。它不具有
$\alpha\Z$ 的形式：否则 $1 = p\alpha$ 且 $\sqrt 2 = q\alpha$
（$p, q \in \Z$）将得 $\sqrt 2 = \frac qp \in \Q$，
矛盾。由上述二分，$G$ 在 $\R$ 中稠密。
\end{solution}

\begin{solution}{exo:b1:reals:10}
设 $s = \sup A > 0$，$t = \sup B > 0$。对 $a \in A$，$b \in B$，有 $ab
\leq st$（\emph{正}数间的不等式可以相乘）。
对 $0 < \varepsilon < \min(s, t)$，取 $a > s - \varepsilon$ 与
$b > t - \varepsilon$，则
\[
ab > (s - \varepsilon)(t - \varepsilon)
= st - \varepsilon(s + t) + \varepsilon^2
> st - \varepsilon (s + t),
\]
而 $\varepsilon(s+t)$ 可任意小，由
上确界的 $\varepsilon$ 刻画（形式为：没有数 $< st$ 是
$AB$ 的上界），故 $\sup AB = st$。

正性条件不可缺少：取 $A = B = \{-1, 0\}$，则 $AB = \{0, 1\}$，
$\sup AB = 1$，而 $\sup A \cdot \sup B = 0 \times 0 = 0$。
\end{solution}

\begin{solution}{exo:b1:reals:11}
记 $s = \sup A$，$i = \inf A$。对 $a, a' \in A$：$a \leq s$
及 $a' \geq i$，得 $a - a' \leq s - i$；由对称性，$\abs{a -
a'} \leq s - i$，故 $s - i$ 是这些差的绝对值的一个上界。对
$\varepsilon > 0$，取 $a > s - \frac\varepsilon2$ 及 $a' < i
+ \frac\varepsilon2$（见\cref{prop:b1:reals:epsilon}及其
对下确界的对偶形式）。于是 $\abs{a - a'} \geq a - a' > s - i -
\varepsilon$。由 $\varepsilon$ 的刻画，
$\operatorname{diam} A = s - i$。

对每个 $a \in A$，有 $i \leq a \leq s$，故 $A \subseteq
\intcc{i}{s}$，是长度为 $\operatorname{diam}
A$ 的一个闭区间。若闭区间 $\intcc{u}{v}$ 包含 $A$，则 $v$ 为
上界，$u$ 为 $A$ 的下界，故 $u \leq i$ 且 $v
\geq s$，从而 $\intcc{i}{s} \subseteq \intcc{u}{v}$。因此
$\intcc{i}{s}$ 即为最小者。
\end{solution}

\begin{solution}{exo:b1:reals:12}
$0 \in E$，故 $E \neq \emptyset$。若 $x > \max(1, y)$，则 $x^2 >
x > y$，故 $E$ 以 $\max(1, y)$ 为上界：$s = \sup E$
存在（见\cref{thm:b1:reals:sup}），且 $s \geq \min(1, y) > 0$
因为 $\min(1, y) \in E$：事实上，若 $y \geq 1$，则 $1^2 = 1
\leq y$，若 $y < 1$，则 $y^2 < y$。

\emph{$s^2 < y$ 不可能。} 取 $0 < h < 1$，使 $h <
\frac{y - s^2}{2s + 1}$。于是
\[
(s + h)^2 = s^2 + 2sh + h^2 \leq s^2 + (2s + 1)h < y ,
\]
故 $s + h \in E$，与 $s$ 为 $E$ 的上界矛盾。

\emph{$s^2 > y$ 不可能。} 取 $0 < h < s$，使 $h <
\frac{s^2 - y}{2s}$。于是 $(s - h)^2 = s^2 - 2sh + h^2 > s^2 -
2sh > y$；对每个 $x \in E$，有 $x^2 \leq y < (s - h)^2$，
从而 $x < s - h$（二者均 $\geq 0$）：$s - h$ 为
$E$ 的一个上界且小于 $s$，与最小性矛盾。

故 $s^2 = y$。唯一性：若 $0 < s < s'$，则 $s^2 <
s'^2$，故两个不同的正根不可能同时平方得 $y$。
单调性：若 $0 < y < y'$，则 $\sqrt y \neq \sqrt{y'}$，且
若 $\sqrt y > \sqrt{y'}$，则平方后得 $y > y'$，故 $\sqrt
y < \sqrt{y'}$。
\end{solution}

\begin{solution}{pb:b1:reals:1}
\textbf{1.} 对 $\frac 18$：$10 = 8 \cdot 1 + 2$，$20 = 8 \cdot 2
+ 4$，$40 = 8 \cdot 5 + 0$；数字为 $1, 2, 5$，余数为 $0$，此后
全为零：$\frac 18 = 0.125$。对 $\frac 17$：$10 = 7 \cdot 1 +
3$，$30 = 7 \cdot 4 + 2$，$20 = 7 \cdot 2 + 6$，$60 = 7 \cdot 8 +
4$，$40 = 7 \cdot 5 + 5$，$50 = 7 \cdot 7 + 1$：数字为 $1, 4, 2,
8, 5, 7$，余数为 $3, 2, 6, 4, 5, 1$。此时余数
已回到 $r = 1$，故六步运算将原样无限重复：
$\frac 17 = 0.(142857)$，余数依次循环取 $1, 3, 2, 6,
4, 5$。

\textbf{2.} $\frac 13$ 以 $2$ 为底（$r_0 = 1$）：$2 = 3 \cdot 0 +
2$，$4 = 3 \cdot 1 + 1$，且 $r = 1$ 重复出现：$\frac 13 =
(0.\overline{01})_2$。$\frac{5}{16}$ 以 $2$ 为底：$10 = 16 \cdot
0 + 10$，$20 = 16 \cdot 1 + 4$，$8 = 16 \cdot 0 + 8$，$16 = 16
\cdot 1 + 0$：$\frac{5}{16} = (0.0101)_2$，有限展开。$\frac 12$
以 $3$ 为底：$3 = 2 \cdot 1 + 1$，且 $r = 1$ 立即重复出现：
$\frac 12 = (0.\overline{1})_3$。

\textbf{3.} 算法在 $N$ 步之后的余数为 $r_N =
b^N p \bmod q$（严格证明见第 14 题，此处只需注意
每一步都把余数乘以 $b$，并
对 $q$ 取余）。之后各位数字全为 $0$，当且仅当某个
$r_N = 0$，即 $q \mid b^N p$；由 $\gcd(p, q) = 1$，
依高斯引理，$q \mid b^N$。若 $q \mid b^N$，则
$q$ 的每一质因数整除 $b^N$，从而整除 $b$（由质数性质）。反之，若
$q = p_1^{a_1} \cdots p_r^{a_r}$ 的每个质因数都整除 $b$，则
取 $A = \max_i a_i$，则每个 $p_i^{a_i}$ 整除 $b^A$，且
$p_i^{a_i}$ 两两互质，故 $q \mid b^A$。对 $q = 20 =
2^2 \cdot 5$：两个质数均整除 $10$（$\frac{1}{20} = 0.05$），
但 $2 \nmid 3$，故 $\frac{1}{20}$ 以 $3$ 为底时无限循环。

\textbf{4.} 由 $1^2 = 1 < 2 < 4 = 2^2$ 得 $s_0 = 1$。于是 $1.4^2
= 1.96 < 2 < 2.25 = 1.5^2$：$\lfloor 10\sqrt 2 \rfloor = 14$，
$s_1 = 1.4$。接着 $1.41^2 = 1.9881 < 2 < 2.0164 = 1.42^2$：$s_2 =
1.41$；$1.414^2 = 1.999396 < 2 < 2.002225 = 1.415^2$：$s_3 =
1.414$；$1.4142^2 = 1.99996164 < 2 < 2.00024449 = 1.4143^2$：$s_4
= 1.4142$。在每种情形，所列不等式恰好说明
$s_n \leq \sqrt 2 < s_n + 10^{-n}$，这正是
$10^n \sqrt 2$ 取整函数的定义。

\textbf{5.} $A_0 = \lfloor x \rfloor = 0$，因 $0 \leq x < 1$。
由 $A_{n-1} \leq b^{n-1} x < A_{n-1} + 1$，两边乘以 $b$：
\[
b\,A_{n-1} \leq b^n x < b\,A_{n-1} + b .
\]
整数 $b\,A_{n-1}$ 满足 $\leq b^n x$，故 $b\,A_{n-1} \leq
A_n$；又 $b^n x < b\,A_{n-1} + b$，而 $b\,A_{n-1} + b$ 为
整数，迫使 $A_n \leq b\,A_{n-1} + b - 1$。因此 $0 \leq d_n =
A_n - b\,A_{n-1} \leq b - 1$：即为一个数字。

\textbf{6.} 裂项相消：$d_k b^{-k} = A_k b^{-k} - A_{k-1}
b^{-(k-1)}$，故
\[
\sum_{k=1}^n d_k b^{-k} = A_n b^{-n} - A_0 = s_n .
\]
将 $A_n \leq b^n x < A_n + 1$ 两边除以 $b^n$，得 $s_n \leq x <
s_n + b^{-n}$。

\textbf{7.} 归纳法：$b^0 = 1 \geq 1$，且 $b^{n+1} = b \cdot
b^n \geq 2(n + 1) \geq n + 2$。单调性：$s_n - s_{n-1} = d_n
b^{-n} \geq 0$。每个 $s_n \leq x$（第 6 题）：$x$ 为
$\{s_n\}$ 的上界。对 $\varepsilon > 0$，由阿基米德性质
可取 $n$ 使 $n + 1 > \frac1\varepsilon$，从而
$b^{-n} < \varepsilon$，于是 $s_n > x - b^{-n} > x -
\varepsilon$（由第 6 题）。由\cref{prop:b1:reals:epsilon}可知 $x =
\sup_n s_n$。

\textbf{8.} 设 $d_k = b - 1$ 对一切 $k > N$ 成立。当 $n > N$ 时，
由有限等比数列求和得
\[
s_n = s_N + (b - 1)\sum_{k=N+1}^{n} b^{-k}
= s_N + b^{-N} - b^{-n} .
\]
故 $x \geq s_n = s_N + b^{-N} - b^{-n}$ 对每个 $n$ 成立；令
末项小于任意 $\varepsilon$（第 7 题），得 $x
\geq s_N + b^{-N}$。但第 6 题在第 $N$ 位处给出 $x < s_N +
b^{-N}$，矛盾。串 $(d_n)$ 为规范串。

\textbf{9.} 有界性：$t_n \leq (b-1)\sum_{k=1}^n b^{-k} = 1 -
b^{-n} < 1$，且 $(t_n)$ 单调不减，故 $y = \sup t_n$
存在，且 $0 \leq y \leq 1$。固定 $n$。对双侧
估计：$t_n \leq y$ 显然。由规范性取 $m > n$，使
$e_m \leq b - 2$。对 $p \geq m$：
\[
t_p - t_n = \sum_{k=n+1}^{p} e_k b^{-k}
\leq (b^{-n} - b^{-p}) - b^{-m} < b^{-n} - b^{-m},
\]
中间和至少损失 $b^{-m}$，相对于全取 $(b-1)$ 的
最大值；当 $p \leq m$ 时，$t_p \leq t_m \leq t_n + b^{-n} -
b^{-m}$ 同样成立（由单调性及 $p = m$ 的情形）。于是每个
$t_p \leq t_n + b^{-n} - b^{-m}$，故 $y \leq t_n + b^{-n} -
b^{-m} < t_n + b^{-n}$。（取 $n = 0$：$y < 1$，故 $y \in
\intco{0}{1}$。）现 $b^n t_n = \sum_{k \leq n} e_k b^{n-k}$ 为
整数，且 $b^n t_n \leq b^n y < b^n t_n + 1$：故 $\lfloor b^n
y \rfloor = b^n t_n$。最后考察 $y$ 的数字：$d_n(y) = b^n
t_n - b \cdot b^{n-1} t_{n-1} = b^n(t_n - t_{n-1}) = e_n$.

\textbf{10.} 第 9 题说：串的值的数字即为该串；
第 5--8 题说：（$x$ 的数字）构成规范串，其
截断的上确界为 $x$（第 7 题）。因此两映射
按两种顺序复合都是恒等映射，故二者互为逆映射，
是 $\intco{0}{1}$ 与规范串之间的双射。若两规范串
的值相等，则对二者施行数字映射将使二者相同：
此即唯一性。这便是\emph{$b$-进
展开定理}。

\textbf{11.} 设 $e_n = b - 1$ 对 $n > M$ 成立，其中 $M \geq 0$ 取最小。
与第 8 题相同，$t_n = t_M + b^{-M} - b^{-n}$ 对 $n \geq M$ 成立，
故其值为 $\sup t_n = t_M + b^{-M}$。若 $M = 0$，
则其值为 $0 + 1 = 1$：以 $10$ 为底时，$0.999\dots = 1$ 精确成立 --- 而
非近似。若 $M \geq 1$，由最小性得 $e_M \leq b - 2$，
于是其值为
\[
t_M + b^{-M} = \frac{b^M t_M + 1}{b^M} \in \intoo{0}{1},
\]
这是一个 $b$-进分数，其\emph{规范}展开为 $e_1 \dots
e_{M-1}\,(e_M + 1)\,000\dots$（有限终止串为规范串，
且其值为同一个数）。反之，有两种表示的实数
必有一种为非规范表示（规范性确定唯一表示，
第 10 题），故必为上述形式。每个
$m/b^N \in \intoo{0}{1}$，其末个非零数字为 $d_N$，
都另有一种非规范表示 $d_1 \dots d_{N-1}(d_N -
1)(b-1)(b-1)\dots$。恰好 $b$-进分数有两种
表示，其余实数都只有一种。

\textbf{12.} 设两串直到第 $m - 1$ 位相同，公共截断为
$P = s_{m-1}$，且 $d_m < e_m$。由第 9 题
（第 $m$ 位处的严格上界估计）：$x < P + d_m b^{-m} + b^{-m} = P +
(d_m + 1)b^{-m} \leq P + e_m b^{-m} \leq y$，末步
依据是 $P + e_m b^{-m}$ 即 $y$ 的截断 $t_m \leq y$。故
$d_m < e_m \implies x < y$；交换角色，$e_m < d_m \implies
y < x$；又两串在第 $m$ 位不同，故二者之一
成立。两方向均得证。

\textbf{13.} 令 $z = bx - A_1 \in \intco{0}{1}$（实际上 $A_1
\leq bx < A_1 + 1$）。对 $n \geq 0$，有 $b^n z = b^{n+1} x - b^n
A_1$，其中 $b^n A_1 \in \Z$，由 $\lfloor u - K \rfloor =
\lfloor u \rfloor - K$（$K$ 为整数），
\[
A_n(z) = A_{n+1}(x) - b^n A_1(x) .
\]
于是 $d_n(z) = A_n(z) - b\,A_{n-1}(z) = A_{n+1} - b^n A_1 -
b\,A_n + b^n A_1 = d_{n+1}(x)$。因此 $bx$ 的小数部分
的数字即为平移后的数字；迭代 $m$ 次，
$b^m x$ 的小数部分的数字为 $(d_{n+m})_{n \geq 1}$。

\textbf{14.} 带余除法：$b^n p = q\,Q_n + r_n$，其中 $0
\leq r_n < q$。以 $q$ 除之得 $b^n x = Q_n + \frac{r_n}{q}$，其中
$0 \leq \frac{r_n}{q} < 1$，故 $Q_n = \lfloor b^n x \rfloor =
A_n$，从而 $A_n = \frac{b^n p - r_n}{q}$。关于递推式：
$b^n p = b(q\,A_{n-1} + r_{n-1}) = q\,(b\,A_{n-1}) +
b\,r_{n-1}$，故 $b^n p$ 与 $b\,r_{n-1}$ 相差
$q$ 的倍数，从而 $r_n = (b\,r_{n-1}) \bmod q$。

\textbf{15.} 将 $b\,r_{n-1}$ 除以 $q$：$b\,r_{n-1} = q\,c +
r_n$，其中 $c = \lfloor b\,r_{n-1}/q \rfloor$。代入
第 14 题中的式子得 $b^n p = q(b\,A_{n-1} + c) + r_n$，
由带余除法的唯一性，得 $A_n = b\,A_{n-1}
+ c$，即 $d_n = c = \lfloor b\,r_{n-1}/q \rfloor$。
第 $n$ 位数字仅依赖于 $r_{n-1}$ --- 这正是第
I 部分的长除法循环，现已得证。

\textbf{16.} 共 $q + 1$ 个余数 $r_0, \dots, r_q$，其取值
都在含 $q$ 个元素的集合 $\intint{0}{q-1}$ 中，由抽屉原理
（\cref{cor:b1:counting:pigeonhole}）知其中有两者相同，不妨记为
$r_N = r_{N+T}$，其中 $0 \leq N < N + T \leq q$。由于 $r_n$
确定 $r_{n+1}$（第 14 题），归纳得 $r_{n+T} =
r_n$（对一切 $n \geq N$）；由于 $r_{n-1}$ 确定 $d_n$
（第 15 题），故 $d_{n+T} = d_n$ 对一切 $n \geq N + 1$ 成立。每个
有理数的展开最终是周期的，前周期 $\leq
q$，周期 $\leq q$。

\textbf{17.} $b^T y$ 的小数部分的数字为
$(d_{n+T}) = (d_n)$（由平移引理及纯周期性），与
$y$ 的规范串相同。由第 10 题，二者的值相等：
$b^T y - \lfloor b^T y \rfloor = y$，故 $(b^T - 1)\,y = \lfloor
b^T y \rfloor = A_T \in \N$，且
\[
y = \frac{A_T}{b^T - 1} ,
\]
有理数，分母整除 $b^T - 1$；分子 $A_T$
即以 $b$ 进制表示时数字依次为 $d_1 \dots d_T$ 的整数。验证：
$0.(142857) = \frac{142857}{999999}$，且 $142857 \times 7 =
999999$，故此即 $\frac 17$。

\textbf{18.} 若 $d_{n+T} = d_n$（对 $n > N$），则 $z$ 为 $b^N x$ 的小数部分，
其数字为 $(d_{N+n})_{n\geq1}$（平移引理），
它们是纯周期的；由第 17 题，$z \in \Q$。于是 $b^N
x = A_N + z$ 得 $x = (A_N + z)/b^N \in \Q$。结合第 16 题：
$x$ 为有理数 $\iff$ 展开最终周期。右端
涉及底数，左端则不涉及：某一底数下的周期性
等价于有理性，因而等价于任一底数下的周期性。

\textbf{19.} 设 $x = \frac 1q$，$r_n = b^n \bmod q$。若
$\gcd(b, q) = 1$，则 $r_T = r_0 = 1$ 当且仅当 $b^T \equiv
1 \pmod q$；这样的 $T$ 存在（由抽屉原理得 $b^i \equiv b^j$，
$i < j$，且 $b$ 模 $q$ 可逆，故 $b^{j-i} \equiv
1$），取最小者——即乘法阶——使余数，
从而数字，呈周期为 $T$ 的纯周期性。
不存在更小的周期：周期 $T'$ 将给出 $(b^{T'}
- 1)\frac1q \in \N$（第 17 题），即 $q \mid b^{T'} - 1$。
取 $q = 7$，$b = 10$：$10 \equiv 3$，$10^2 \equiv 2$，$10^3
\equiv 6$，$10^4 \equiv 4$，$10^5 \equiv 5$，$10^6 \equiv 1
\pmod 7$：阶为 $6$，且 $\frac 17$ 的周期确为六。

\textbf{20.} 该串含无穷多个 $0$（$x^*$ 的数字多数为零），
故它是规范的，且 $x^*$ 有良定义
（第 9 题）。假设数字最终呈周期性，
周期为 $T$，且自第 $N$ 位之后成立。无穷多个数字等于 $1$（每个
三角数处一个），故某个 $1$ 位于 $j > N$ 处；于是
周期性使 $1$ 出现在每个位置 $j + kT$ 上：自 $j$
起，相邻 $1$ 之间的间隔至多为 $T$。但
$1$ 恰位于三角数处，其相邻间隔
为 $\frac{(j+1)(j+2)}{2} - \frac{j(j+1)}{2} = j + 1$，最终超过 $T$：
矛盾。故数字最终不呈周期性，由第 18 题，
$x^* \notin \Q$ --- 无理性
仅由数字的节奏即可读出。

\textbf{21.} 给定 $x < y$，由第 7 题取 $n$ 使 $b^{-n}
< y - x$；令 $m = \lfloor b^n x \rfloor + 1$。则 $b^n x < m
\leq b^n x + 1 < b^n y$，于是 $x < \frac{m}{b^n} < y$：稠密性，
对任一底数同时成立（取 $b = 2$ 即得
\cref{exo:b1:reals:8}）。对 $\frac pq \in \intoo{0}{1}$，取底数
$b = q$：首位数字为 $\lfloor q \cdot \frac pq \rfloor =
p$，而 $q \cdot \frac pq = p$ 的小数部分为 $0$：此后
各位数字全为零，是有限展开 $\frac pq =
(0.p)_q$。有限展开依赖于底数；而周期性 ---
即有理性 --- 则不依赖底数（第 18 题）。

\textbf{22.} 每个 $e_k \in \{1, 2\}$ 都是底数 $10$ 的数字，且
该串永不以全为 $9$ 的尾部结束，是规范的。其值 $y$ 位于
$\intco{0}{1}$ 中，其数字恰为 $(e_k)$（第 9 题）。固定
$k$：第 $k$ 位上 $y$ 的数字为 $e_k$，它 $\neq$ 第 $k$ 位
$x_k$ 的数字，故 $y$ 与 $x_k$ 的规范串不同，从而
$y \neq x_k$（第 10 题：编码是单射）。因此 $y$ 不在
任何列表中：映射 $\N^* \to \intco{0}{1}$ 皆非满射。实数
与有理数不同，无法被枚举 ---
不可数性，其理论在\cref{ch:b1:topology}中展开。

\textbf{23.} 若两展开式至第 $n$ 位一致，则 $x$ 与 $y$
有相同的截断 $s_n$，由第 6 题，二者均落在
$\intco{s_n}{s_n + b^{-n}}$ 中，这是长为 $b^{-n}$ 的区间：
$\abs{x - y} < b^{-n}$。反过来：$x = 0.1$ 与 $y = 0.0999$
（有限展开，因而规范）满足 $\abs{x - y} = 10^{-4} <
10^{-3}$，但二者的展开式在第一位即不同。
症结在于第 11 题中的 $b$ 进分数：在它们附近，
微小的移动便会改变所有显示的数字（$0.0999 \to
0.1000$），因为它们恰好是那些
非规范孪生串潜伏其中的实数。

\textbf{24.} 取 $p = 1$，$q = 10$，$b = 2$，$r_0 = 1$：$2 = 10 \cdot
0 + 2$，$4 = 10 \cdot 0 + 4$，$8 = 10 \cdot 0 + 8$，$16 = 10
\cdot 1 + 6$，$12 = 10 \cdot 1 + 2$ --- 且 $r_5 = 2 = r_1$：
余数以 $(2, 4, 8, 6)$ 为循环，自下标 $1$ 起。数字：$d_1 =
0$，随后的循环节 $d_2 d_3 d_4 d_5 = 0, 0, 1, 1$：
\[
\tfrac{1}{10} = (0.0\overline{0011})_2 ,
\]
前周期 $1$，周期 $4$。由第 3 题，有限的 $2$ 进制展开
要求 $10$ 的每个质因数都整除 $2$；质数
$5$ 不满足。因此 $0.1$ \emph{不能}
由任何有限的二进制串表示 --- 存储有限位数的计算机
只保留一个截断，累积的截断误差正是
浮点数 $0.1 + 0.2$ 与 $0.3$ 在末几位不同的原因。

\textbf{25.} (i) 完备性产生了这些值：$x = \sup s_n$
与 $y = \sup t_n$（第 7、9 题）--- 若仅在 $\Q$ 内，
$\sqrt 2$ 的规范串将一无所指。（ii）阿基米德性质使
$b^{-n}$ 最终小于任一
$\varepsilon$，迫使截断逼近其
上确界（第 7、21 题）。（iii）取整函数的唯一性条款
确定了第 14 题中的 $Q_n = A_n$，并使每一次
数字提取 $\lfloor u - K \rfloor = \lfloor u \rfloor - K$ 合法
（第 13 题）。（iv）抽屉原理应用于有限多个
余数，是周期性的唯一动力（第 16 题）。
寓意：规范串忠实地编码 $\intco{0}{1}$，并将
有理性转化为可见的节奏；然而数字串的加法
需要进位从无穷远的右侧传递而来，因此没有
有限步的规则甚至连和的首位数字都算不出 ——
而\cref{thm:b1:reals:sup}所给的上确界接口
以一条公理处理全部分析。数字是一幅绝妙的
\emph{图景}，描绘的是 $\R$；上确界则是其\emph{引擎}。
\end{solution}
